% !TEX root = ../main_zh.tex
\chapter{货币数量论与通货膨胀}
\label{ch:quantity}

第~\ref{ch:money}~章讲述了货币是如何被创造出来的——中央银行与商业银行如何携手扩张在经济中流通的通货与存款。本章顺势追问下一个问题：当这一存量增长时，价格水平会发生什么？简短的回答，也是古典货币经济学的核心思想，是：价格水平归根结底是一种货币现象。当货币数量的增长快于商品数量时，每一单位货币所追逐的商品比从前更多，价格便随之攀升。价格水平持续而广泛的上涨，就是我们所说的\emph{通货膨胀}（inflation），而在长期内，它由货币的增长率所驱动。

在建立把这一点说清楚的模型之前，我们先停下来讨论一个学宏观经济学的人几乎总是搞反的前置问题：\emph{通货膨胀究竟为什么重要？}流行的抱怨是，通货膨胀侵蚀我们的购买力，让我们变穷。我们将看到，这种直觉若照字面理解，在很大程度上是一种混淆——在一个无摩擦的世界里，收入会随着价格一同上涨，通货膨胀真正的实际成本另有所在。弄清这些成本究竟在哪里，以及哪里甚至还存在某些抵消性的好处，能让我们更清楚地理解货币政策能做什么、不能做什么。随后本章发展出\emph{货币数量论}（quantity theory of money），它由方程 $MV = PY$ 概括；由此推导出通货膨胀率的一个简洁表达式；再把同一方程改写成一种货币需求理论；最后以\emph{古典二分法}（classical dichotomy）与\emph{货币中性}（neutrality of money）收尾——后者断言，货币在长期内不影响实际量。这一断言只有在价格已充分调整时才成立；中性在短期（价格具有粘性时）的失效，恰恰是通向后续各章凯恩斯模型的入口。

\section{作为一种税的通货膨胀}

观察通货膨胀的一个有用的初步视角是财政视角。当政府无法或不愿通过普通的征税或借债来为其支出筹资时，它可以干脆印钞来付账。它每发行一单位新货币，就稀释了已经在流通中的每一单位货币的购买力。货币的持有者由此在实际意义上变穷了，而这些货币所支配的资源被转移给了政府。其效果与对货币余额征收的一种税无异。

\defn{铸币税（通胀税）}{
\emph{铸币税}（seigniorage）是政府通过发行货币获得的实际收入。等价地，它就是\emph{通胀税}（inflation tax）：当货币存量增长、价格水平上升时，公众持有的货币的实际价值下降，而这部分失去的购买力归于货币的发行者。这里的“税基”是公众选择持有的实际货币余额 $M/P$，“税率”则是侵蚀它们的通货膨胀率 $\pi$。
}

这就是认为通胀相当于在征收通胀税的传统观点。它值得记住，因为它解释了为什么财政承压的政府如此频繁地诉诸印钞机，也解释了为什么恶性通货膨胀（hyperinflation）的发作几乎总是财政崩溃的发作：通胀税是最后的收入来源。它还把货币数量论的核心问题用财政语言框定下来——如果政府不断扩张 $M$ 来为支出筹资，这对 $P$ 会有什么影响？

\section{通货膨胀的成本}

如果通货膨胀转移了资源、扭曲了决策，那么这些成本究竟由谁、在何处承担？把经济理论所识别出的真正一阶成本，与公众最为恐惧、但实际上更多是表象而非实质的那种成本区分开来，是有益的。

\subsection{真实的成本}

\begin{enumerate}
    \item \textbf{鞋底成本（交易成本上升）。}当预期价格会持续上涨时，持有货币的代价变高，于是人们设法节约自己的货币余额——更频繁地往返银行、把财富存放于带息资产或实物资产、缩短合同的期限。长期合同变得更难订立，因为双方都必须不断把变动的价格水平考虑进来。这个名称让人联想到额外跑银行磨损的鞋底；更广义地说，它涵盖了在价格稳定时本无必要、却因管理货币持有而耗费的全部实际资源。

    \item \textbf{菜单成本（价格调整的成本）。}企业必须实实在在地更改其标价——重印菜单和目录、重贴货架标签、重新设定系统——而每一次这样的改动都消耗实际资源。通货膨胀率越高，价格就必须越频繁地修订，这些成本也就越大。

    \item \textbf{资源配置被扭曲。}价格是引导资源流向其最有价值用途的信号。当总体价格水平在上涨、且不同商品以不同速度上涨时，就难以分辨某种商品\emph{相对}价格的变化（这本应重新引导资源）与总体价格水平的变化（这本不该引导资源）。\emph{价格的信号作用失效}，资源因此被错误配置。

    \item \textbf{债务人与债权人之间的再分配。}当通货膨胀高于签订贷款时所预期的水平时，实际支付的\emph{实际}利率被压低。债务人用价值低于预期的货币偿还，而债权人收到的实际价值少于合同所承诺的。因此通货膨胀把财富从债权人（贷方）再分配给债务人（借方）：\emph{借方得益、贷方受损}。
\end{enumerate}

\subsection{为什么“通货膨胀缩水了我的收入”不是一阶成本}

请注意上面这份清单里\emph{缺}了什么：大多数人恐惧通货膨胀的理由，也就是价格上涨缩水了他们收入的购买力、从而让他们变穷。这一顾虑在理论上站不住脚。如果经济中所有价格按比例同步上涨，那么工资、租金和利润——它们本身也是价格——就理应同比例上涨。一个工资和食品账单都翻了一倍的工人，处境并没有变差。实际收入，即名义收入与价格水平之比，保持不变。统一且被充分预期到的通货膨胀，本身并不会减少任何人的实际收入。

\rmkb[流行的担忧在何时变成真问题]{上面这个理论论证假定所有价格和收入都同步、即时地一起变动。现实在两方面偏离了这一假定。第一，通货膨胀往往是\emph{结构性}的：不同价格以不同速度上涨，因此某个特定家庭所购买的篮子的涨速，可能快于它那一份特定的收入来源。第二，工资和许多价格具有\emph{粘性}——它们不会立即调整。当商品价格上涨、而某工人的名义工资被合同锁定一年时，这位工人的实际收入在这段期间里确实下降了。正是这些摩擦，让流行的担忧在短期里有了着力点。它们也是凯恩斯主张的基础，即财政政策和货币政策在短期内能够推动实际产出：如果价格能即时调整，就没有可供利用的短期实际效应了。我们将在本章末尾讨论货币中性时回到这一点，并在第~\ref{ch:keynes}~章至第~\ref{ch:adas}~章中将其充分展开。}

\section{通货膨胀的收益}

通货膨胀并非全是成本。适度的通货膨胀率带来两项值得明确指出的好处，二者都经由上面那些成本所通过的同一渠道发挥作用。

\begin{enumerate}
    \item \textbf{它刺激了信贷需求。}通过压低存量债务的实际利率，通货膨胀使借贷在实际意义上变得更便宜，从而鼓励企业和家庭举债与投资。更高的信贷需求能够支撑支出与经济活动。

    \item \textbf{它增加了劳动力市场的弹性。}名义工资是出了名地难以下调：工人抵制直接的减薪，合同与惯例使其\emph{向下刚性}。然而劳动力市场有时需要实际工资下降——例如当一家企业面临需求疲软、否则就不得不裁员时。通货膨胀提供了一条出路。如果价格水平上涨而名义工资保持不变，那么\emph{实际}工资就在没有任何人被正式减薪的情况下下降了。更低的实际工资使企业雇佣工人的成本下降，从而提高劳动需求、支撑就业。如此一来，适度的通货膨胀就让劳动力市场运转得更顺畅。
\end{enumerate}

\section{货币数量论}

我们现在从“通货膨胀为什么重要”转向“是什么决定了通货膨胀”。古典的回答是\emph{货币数量论}，它从一个把货币数量与它所完成的交易价值联系起来的会计恒等式出发。

\defn{数量方程}{
\emph{数量方程}（quantity equation，又称交换方程）指出，货币数量乘以它的流通速度，等于交易的货币价值：
\[
M \times V = P \times T,
\]
其中 $M$ 是名义货币存量，$V$ 是\emph{货币的流通速度}（velocity of money，即每单位货币在一个时期内平均转手的次数），$P$ 是价格水平，$T$ 是交易量。
}

交易量 $T$ 很难直接衡量，所以在实践中我们用实际产出（实际 GDP）来替代它，记为 $Y$。经此替换，右端的 $P \times Y$ 恰好就是\emph{名义} GDP，于是数量方程就成为我们全书将要使用的形式：
\begin{equation}
M \times V = P \times Y.
\label{eq:qty-eqn}
\end{equation}

把这个恒等式转化为一种\emph{理论}的，是一组关于哪些项可以自由变动、哪些项在短期内固定的陈述。

\asm{货币数量论的假设}{
\begin{description}
    \item[产出由供给侧决定。]实际产出 $Y$ 由经济的潜力——它的资本、劳动和技术——所决定，因此在短期内被视为给定，与货币存量无关。
    \item[流通速度由习惯决定。]流通速度 $V$ 由经济的交易技术与支付习惯（人们如何被支付、多久支付一次、以及如何持有货币）所支配。这些习惯只会缓慢改变，因此 $V$ 在短期内大致不变。
\end{description}
}

\rmkb[流通速度与利率]{流通速度并非字面意义上的常数。当持有货币的机会成本上升时它会上升：当利率攀升——也就是流动性的价格上涨——时，人们相对于其支出会持有更少的货币，于是每单位货币转手得更频繁，$V$ 随之增加。我们的假设仅仅是，在短期内，$V$ 相较于我们所关心的 $M$ 与 $P$ 的波动而言变动甚微。$V$ 对利率的这种依赖，正是第~\ref{sec:money-demand}~节的流动性偏好改写所要明确刻画的。}

把 $V$ 与 $Y$ 固定为背景，方程~\eqref{eq:qty-eqn} 就把价格水平直接同货币存量绑在了一起。要干净利落地看清这一点，对两边取对数，
\[
\log M + \log V = \log P + \log Y,
\]
再对时间求导。由于 $\log X$ 的时间导数就是增长率 $\Delta X / X$，这便得到
\begin{equation}
\frac{\Delta M}{M} + \frac{\Delta V}{V} = \frac{\Delta P}{P} + \frac{\Delta Y}{Y},
\label{eq:qty-growth}
\end{equation}
或用百分比变化的语言来写，
\[
\%\Delta M + \%\Delta V = \%\Delta P + \%\Delta Y.
\]
方程~\eqref{eq:qty-growth} 把数量方程表达为一种\emph{增长率}之间的关系：货币增长率与流通速度增长率之和，等于通货膨胀率与实际产出增长率之和。由此可得两条结论。

\state{货币数量论的两条结论}{
\begin{enumerate}
    \item \textbf{通货膨胀由货币增长所驱动。}在流通速度与实际产出的增长率固定不变时，货币增长率 $\Delta M / M$ 的任何上升都会一对一地传递到通货膨胀率 $\Delta P / P$。持续的通货膨胀，归根结底是持续的货币创造的后果。
    \item \textbf{货币增长追随名义 GDP 增长。}由于 $\Delta V / V \approx 0$，货币增长 $\Delta M / M$ 与名义 GDP 的增长 $\%\Delta P + \%\Delta Y$ 大致同步变动。在现实中，中央银行扩张货币供给的速度，大体上与名义产出的增长相匹配。
\end{enumerate}
}

在~\eqref{eq:qty-growth} 中令流通速度增长 $\Delta V / V$ 为零，并回忆通货膨胀率 $\pi$ 就是价格水平的增长率，$\pi = \Delta P / P$，我们就能直接解出通货膨胀：
\begin{equation}
\pi = \frac{\Delta M}{M} - \frac{\Delta Y}{Y}.
\label{eq:inflation-rate}
\end{equation}

\defn{由货币增长决定的通货膨胀率}{
在流通速度不变的条件下，通货膨胀率等于货币增长与实际产出增长之差，
\[
\pi = \frac{\Delta M}{M} - \frac{\Delta Y}{Y}.
\]
}

方程~\eqref{eq:inflation-rate} 给出了对通货膨胀一个粗略却富有启发的估计。可以这样理解它：理论上，流通中的货币数量“理应”与经济的实际产出同步增长，使得货币恰好够在价格不变的情况下完成那些新增的实际交易。货币增长中\emph{超出}实际增长的那一部分，没有额外的商品可供追逐，于是溢出为更高的价格。这一超额的部分，就是通货膨胀率。

\section{货币需求与流动性偏好}
\label{sec:money-demand}

到目前为止，我们一直把数量方程读作一种价格水平理论。同一方程经过重排，也是一种\emph{货币需求}理论——关于人们愿意持有多少货币。把 $MV = PY$ 重排以分离出实际余额，
\begin{equation}
\frac{M}{P} = \frac{Y}{V}.
\label{eq:real-money}
\end{equation}
左端的 $M/P$ 是\emph{实际货币供给}——货币存量的购买力——在均衡时它等于公众所需求的实际余额数量。因此数量方程刻画了\emph{流动性偏好}：维持经济体运转所需要的货币量。由此可得两点观察。第一，这一方程从根本上讲的是货币\emph{需求}。第二，它刻画的是一种\emph{均衡}：当货币市场出清时，方程中的 $M$ 同时代表货币的供给与对货币的需求。

凯恩斯把~\eqref{eq:real-money} 的右端推广为一个明确的货币需求函数。他不再以固定的流通速度把实际余额钉在 $Y/V$ 上，而是让对实际余额的需求通过一个\emph{流动性偏好}（liquidity preference）函数 $L$ 依赖于收入与利率：
\begin{equation}
\pr{\frac{M}{P}}^{\!d} = L(i, Y) = L\pr{r + \pi^{e},\, Y}.
\label{eq:liq-pref}
\end{equation}
这里 $i$ 是名义利率，由费雪关系（Fisher relation）它等于实际利率 $r$ 与预期通货膨胀 $\pi^{e}$ 之和，即 $i = r + \pi^{e}$。两个自变量的符号至关重要：
\begin{enumerate}
    \item \textbf{货币需求随收入 $Y$ 递增。}更高的实际收入水平意味着更多的支出、更多需要融资的交易，因此需要更多的货币来完成它们。于是 $L$ 随 $Y$ 上升。
    \item \textbf{货币需求随名义利率 $i$ 递减。}利率是持有货币的机会成本：以现金或无息账户形式持有的货币，放弃了债券及其他资产所能提供的回报 $i$。当 $i$ 上升时，持有货币的代价变高，于是人们持有得更少——等价地说，流通速度上升，每单位货币承担更多的交易。于是 $L$ 随 $i$ 下降。
\end{enumerate}

\rmkb[为什么货币需求用的是\emph{预期}通货膨胀]{对于一项\emph{今天}就要做出的、关于在接下来一个时期内持有多少货币的决策而言，真正重要的机会成本是\emph{名义}利率，而相关的名义利率，是那个为储蓄者\emph{预期}会发生的通货膨胀作补偿的利率，而非事后才实现的通货膨胀。货币需求是一个\emph{事前}（ex ante，即前瞻性）的概念，所以它依赖于\emph{预期}通货膨胀 $\pi^{e}$，这正是为什么 $L$ 的第二个自变量写作 $r + \pi^{e}$ 而非 $r + \pi$。}

在市场的另一侧，名义货币供给 $M$ 由中央银行设定。我们把它视为一个外生的、独立的政策决策：货币当局通过行政手段选择 $M$，而这一选择被模型的其余部分当作给定。货币市场的均衡要求实际货币供给等于对它的实际需求，
\begin{equation}
\frac{M}{P} = L(i, Y) = L\pr{r + \pi^{e},\, Y}.
\label{eq:money-eqm}
\end{equation}
在名义货币供给 $M$ 由政策固定的情况下，利率 $i$ 与价格水平 $P$ 会进行调整，以使~\eqref{eq:money-eqm} 成立。这一均衡条件就是我们将在第~\ref{ch:islm}~章构建的 LM 关系的引擎。

\section{古典二分法与货币中性}

我们现在可以陈述货币理论的核心长期命题，以及它在短期内可能失效的确切含义——通往后续一切内容的桥梁。

\defn{古典二分法}{
\emph{古典二分法}是把经济变量分离为\emph{实际}量（产出、就业、相对价格、实际利率）与\emph{名义}量（货币存量、价格水平、名义工资）。这一二分法认为，名义变量不会影响实际变量：货币数量的变化推动价格水平，却不触动实际配置。
}

如果古典二分法成立，那么价格水平就不能改变实际经济行为。货币存量的变化以及产生它的政策，对产出、就业或消费都没有影响——它们只是把所有名义量一起按比例放大或缩小。这时我们就说货币是\emph{中性}的。

\thm{货币中性}{
如果名义货币存量的变化使所有名义变量按同一比例变化、而让所有实际变量保持不变，那么货币就是\emph{中性}的。当经济中\emph{每一个}价格都在同一时间、朝同一方向、按同一比例进行调整，从而没有任何相对价格被扰动时，中性成立。
}

中性成立的条件是苛刻的：它要求所有价格——商品价格、工资、资产价格——都即时且按比例地一起变动。当这一条件被满足时，货币存量翻倍只是简单地使价格水平以及每一项名义工资和合同翻倍，而实际工资、实际产出和实际利率都原封不动。这与第~\ref{ch:money}~章简单货币模型中遇到的中性结果是同一个，如今则在整个经济的层面上陈述出来。

然而现实充满摩擦。价格与工资具有粘性：它们被提前设定、被不频繁地修订，且并非全都同时变动。当货币存量变化、而价格尚未跟上时，相对价格被暂时扰动，于是货币的变化\emph{确实}推动了实际量。因此在短期内，\emph{货币不是中性的}，货币政策能够影响产出与就业的水平。

\state{中性是长期性质，而非短期性质}{
\begin{itemize}
    \item \textbf{短期：}价格具有粘性，古典二分法失效，货币\emph{不}中性。货币政策与财政政策能够推动实际产出与就业。
    \item \textbf{长期：}只要给予足够的时间，每一个有粘性的价格都会调整，直到经济达到一个不再残留任何价格刚性的新均衡。古典二分法得以恢复，货币是中性的。政策只推动价格水平，而非实际量。
\end{itemize}
}

这一区分是本课程余下部分的组织原则。经济存在一个\emph{长期}均衡，它锚定于潜在产出、由供给侧支配，在那里货币是中性的、古典二分法成立。它也存在\emph{短期}波动——经济周期——在此期间，有粘性的价格让需求侧的扰动（以及更难以控制的供给侧冲击，例如全要素生产率的变动）把经济推离潜在水平。随着时间推移，价格、利率及相关变量的调整，会把一个偏离了潜在产出的短期均衡重新拉回潜在产出。

\rmkb[凯恩斯论短期与长期]{凯恩斯的理论正是建立在短期与长期这一区分之上的，二者以价格是否具有粘性来划分。在短期，供给是可变的、价格是刚性的，于是货币政策与财政政策可以移动总需求、从而改变产出水平；货币不是中性的。在长期，价格完全可变，总供给在潜在产出处保持刚性，货币是中性的。凯恩斯框架的承诺在于：政策能够熨平经济周期的短期波动，尽管它无法改变经济的长期潜力。后续各章的模型——凯恩斯交叉（第~\ref{ch:keynes}~章）、IS--LM（第~\ref{ch:islm}~章）以及总需求与总供给（第~\ref{ch:adas}~章）——就是把这一承诺付诸运转的机器。}
